% chktex-file 3
% chktex-file 8
\section{Artificial de-rendering and re-rendering mechanics}\label{sec:derendering}

The term \emph{de-rendering} denotes a controlled removal of a localized
visible character block from the active boundary-to-bulk projection.  It is
not a destruction of the underlying Hilbert-space state.  Let
\(\mathcal C_{\mathrm{loc}}\subseteq\mathcal C\) be the characters assigned to
the selected region and let
\(\Pi_{\mathrm{loc}}=\sum_{c\in\mathcal C_{\mathrm{loc}}}|c\rangle
\langle c|\).  The two dark ledgers used by the channel are
\begin{equation}
  c_{\mathrm{dark}}^{\mathrm{res}}
  =\frac{834433}{362670}
  =\SimOutputDarkResidual,
  \qquad
  c_{\mathrm{dark}}^{\mathrm{comp}}
  =c_{\mathrm{dark}}^{\mathrm{res}}+1
  =\frac{1197103}{362670}
  =\SimOutputDarkCompleted.
  \label{eq:dark-ledgers-derendering}
\end{equation}
Their ratio
\begin{equation}
  \eta_D
  =\frac{c_{\mathrm{dark}}^{\mathrm{res}}}
  {c_{\mathrm{dark}}^{\mathrm{comp}}}
  =\frac{834433}{1197103}
  \label{eq:derendering-channel-ratio}
\end{equation}
sets the relative residual-sector amplitude.

\subsection{Modular decoupling channel}

A norm-preserving realization of the modular decoupling operator is the
Stinespring map
\begin{equation}
  \widehat{\mathcal D}_{\mathrm{derender}}
  =\sum_{c\in\mathcal C_{\mathrm{loc}}}
  \left(
    \sqrt{\eta_D}\,
    |\chi_{\mathrm{res}}^c,0\rangle
    +\sqrt{1-\eta_D}\,
    |\chi_{\mathrm{comp}}^c,1\rangle
  \right)\langle c|,
  \label{eq:modular-decoupling-operator}
\end{equation}
where the last label is an auxiliary ledger flag and the dark states are
orthonormal for distinct \(c\).  Consequently,
\begin{equation}
  \widehat{\mathcal D}_{\mathrm{derender}}^{\dagger}
  \widehat{\mathcal D}_{\mathrm{derender}}
  =\Pi_{\mathrm{loc}},
  \qquad
  \operatorname{Tr}\rho_{\mathrm{dark}}
  =\operatorname{Tr}(\Pi_{\mathrm{loc}}\rho_{\partial}),
  \label{eq:derendering-isometry}
\end{equation}
with
\begin{equation}
  \rho_{\mathrm{dark}}
  =\widehat{\mathcal D}_{\mathrm{derender}}
  \Pi_{\mathrm{loc}}\rho_{\partial}\Pi_{\mathrm{loc}}
  \widehat{\mathcal D}_{\mathrm{derender}}^{\dagger}.
  \label{eq:dark-state-after-decoupling}
\end{equation}
The orthogonal visible complement is left unchanged.  Thus the complete
channel is trace preserving when Eq.~\eqref{eq:dark-state-after-decoupling}
is combined with
\((1-\Pi_{\mathrm{loc}})\rho_{\partial}(1-\Pi_{\mathrm{loc}})\).

For a local gauge-rendering cost
\begin{equation}
  \mathcal C_{\overline B}(\rho)
  =\gamma_3\operatorname{Tr}(\rho Q_{\mathrm{SU}(3)}^2)
  +\gamma_{\mathrm{EM}}\operatorname{Tr}(\rho Q_{\mathrm{EM}}^2)
  +\gamma_{\mathrm{fr}}\Delta_{\mathrm{fr}}^2,
  \label{eq:gauge-rendering-cost}
\end{equation}
decoupling removes the local visible contribution while preserving its
charge labels in the dark ledger.  Since the branch is unchanged,
\(\Delta_{\mathrm{fr}}=0\) throughout the channel.

\subsection{What is nullified}

Write the Lorentzian projection as a background plus an actively controlled
deformation,
\begin{equation}
  g_{\mu\nu}^{(L)}
  =\eta_{\mu\nu}+\delta g_{\mu\nu}^{\mathrm{act}}[\rho_E].
  \label{eq:active-metric-decomposition}
\end{equation}
The decoupled characters no longer contribute source terms to the local
entropy-flow map, so the decoupling limit is
\begin{equation}
  \delta g_{\mu\nu}^{\mathrm{act}}\longrightarrow0,
  \qquad
  \beta_x\longrightarrow0,
  \qquad
  g_{\mu\nu}^{(L)}\longrightarrow\eta_{\mu\nu}.
  \label{eq:derendering-background-restoration}
\end{equation}
Thus the active metric slice is nullified, but the spacetime metric itself is
not.  A total limit \(g_{\mu\nu}\to0\) is forbidden by the normalized Gram
construction.  In fact,
\begin{equation}
  g^{(G)}
  =\frac{\operatorname{Sym}(\widetilde g+\epsilon_{\mathrm{eig}}I_4)}
  {\operatorname{Tr}[\operatorname{Sym}(\widetilde g+
  \epsilon_{\mathrm{eig}}I_4)]},
  \qquad
  \operatorname{Tr}g^{(G)}=1,
  \label{eq:normalized-gram-derendering}
\end{equation}
and
\begin{equation}
  \lambda_{\min}(g^{(G)})
  \geq
  \frac{\epsilon_{\mathrm{eig}}}
  {\operatorname{Tr}\widetilde g+4\epsilon_{\mathrm{eig}}}
  >0.
  \label{eq:derendering-eigenvalue-floor}
\end{equation}
The zero tensor would instead have trace and determinant equal to zero,
contradicting Eqs.~\eqref{eq:normalized-gram-derendering} and~\eqref{eq:derendering-eigenvalue-floor}.

\subsection{Conserved ledgers}

The de-rendering channel is admissible only when the controller is included
in the closed system.  On a Cauchy surface \(\Sigma\), total four-momentum
obeys
\begin{equation}
  P_{\mathrm{tot}}^{\mu}
  =\int_{\Sigma}d\Sigma_{\nu}
  \left(
    T_{\mathrm{vis}}^{\mu\nu}
    +T_{\mathrm{dark}}^{\mu\nu}
    +T_{\mathrm{ctrl}}^{\mu\nu}
  \right),
  \qquad
  \nabla_{\mu}T_{\mathrm{tot}}^{\mu\nu}=0,
  \qquad
  \Delta P_{\mathrm{tot}}^{\mu}=0.
  \label{eq:derendering-momentum-conservation}
\end{equation}
Charge conservation is expressed by an intertwining relation between the
visible charge and its dark topological encoding,
\begin{equation}
  Q_{\mathrm{dark}}
  \widehat{\mathcal D}_{\mathrm{derender}}
  =\widehat{\mathcal D}_{\mathrm{derender}}Q_{\mathrm{vis}},
  \qquad
  \Delta\langle Q_{\mathrm{tot}}\rangle=0.
  \label{eq:derendering-charge-conservation}
\end{equation}
Finally, the information ledger is bounded by the saturated horizon budget
\begin{equation}
  N_{\mathrm{sat}}
  =\frac{3\pi}{L_P^2\Lambda_{\mathrm{holo}}}
  =\SimOutputSaturatedBitBudget\ \text{bits},
  \qquad
  N_{\mathrm{act}}+N_{\mathrm{dark}}+N_{\mathrm{free}}
  =N_{\mathrm{sat}}.
  \label{eq:saturated-bit-conservation}
\end{equation}
At the \(10\,\mathrm{m}\) benchmark the addressable local allocation is
\begin{equation}
  N_{\mathrm{local}}
  =\SimOutputLocalMemoryBits\ \text{bits},
  \qquad
  N_{\mathrm{free}}^{\mathrm{after}}
  =N_{\mathrm{free}}^{\mathrm{before}}+N_{\mathrm{act}}^{\mathrm{before}}.
  \label{eq:local-bit-release}
\end{equation}

\subsection{Destination reconstruction}

Let \(\widehat T_{\partial}(\bm{x}_{\mathrm{tar}})\) relabel the boundary
address of the stored character block.  The inverse channel followed by
phase-locked excitation defines
\begin{equation}
  \widehat{\mathcal R}_{\mathrm{rerender}}(\bm{x}_{\mathrm{tar}})
  =\widehat{\mathcal O}_{\mathrm{excitation}}(\theta)
  \widehat T_{\partial}(\bm{x}_{\mathrm{tar}})
  \widehat{\mathcal D}_{\mathrm{derender}}^{\dagger},
  \label{eq:rerendering-operator}
\end{equation}
which reconstructs
\begin{equation}
  \rho_{\partial}^{\mathrm{tar}}
  =\widehat{\mathcal R}_{\mathrm{rerender}}
  \rho_{\mathrm{dark}}
  \widehat{\mathcal R}_{\mathrm{rerender}}^{\dagger}
  \longmapsto
  g_{\mu\nu}^{(L)}(\bm{x}_{\mathrm{tar}}).
  \label{eq:rerendered-target-slice}
\end{equation}
This is non-local only as a boundary-address operation: the formal map does
not contain a sequence of intermediate bulk slices.  It is not, by itself, a
proof of acausal transport.  A physical protocol must still satisfy causal
authorization of the target operation,
\begin{equation}
  \bm{x}_{\mathrm{tar}}\in J^+(\bm{x}_{\mathrm{src}}),
  \qquad
  t_{\mathrm{tar}}-t_{\mathrm{src}}
  \geq\frac{d_{\partial}(\bm{x}_{\mathrm{src}},
  \bm{x}_{\mathrm{tar}})}{c},
  \label{eq:rerendering-causal-authorization}
\end{equation}
unless an independently specified causal boundary dynamics proves a stronger
communication channel.